\Hefteintrag{1.3}{Besondere Methoden: Main/Hauptprogramm}{%
Möchte man für ein Java-Programm starten, benötigt man eine \emphColA{\LoesungLuecke{main-Methode}{8cm}}.

Diese nennt man auch den \emphColA{Einstiegspunkt/Hauptprogramm} eines Programms. In ihr werden meistens die benötigten \LoesungLuecke{Objekt}{5cm} erstellt und ihre \LoesungLuecke{Methoden}{5cm} aufgerufen.

Der \emphColC{Kopf} der Main-Methode ist immer gleich (siehe Bsp.):

\emphColC{public static void main() \{}\\
~~~~~~\emph{Spieler s1 = new Spieler(100,50);}\\
~~~~~~\emph{// ... weitere Objekte und Methoden}\\
\emphColC{\}}

Das Schlüsselwort \LoesungLuecke{static}{3cm} bedeutet hierbei, dass kein Objekt benötigt wird, um die Methode aufzurufen. In Blockly2Java wird die main-Methode automatisch zum Ausführen ausgewählt.
}